%% ════════════════════════════════════════════════════════════════════════════
\section{Conclusion}
\label{sec:conclusion}
%% ════════════════════════════════════════════════════════════════════════════

\subsection{Summary of Contributions}
\label{subsec:contributions}

This paper presented the design, implementation, and evaluation of the
Smart Campus Classroom Management System---an end-to-end IoT and edge-AI
platform that addresses three concrete problems in university classroom
management: manual attendance, unmonitored environment, and the absence
of early warning for at-risk students.

The principal contributions are as follows:

\begin{enumerate}
  \item \textbf{A fully self-contained edge stack.}  All services
    (MQTT broker, backend, database, AI inference, frontend) run on a
    single Raspberry Pi~4B without cloud dependency, making the system
    deployable in institutions with restricted internet access or strong
    data-residency requirements.

  \item \textbf{Snapshot-per-cycle attendance model.}  A bidirectional
    present$\leftrightarrow$absent evaluation loop that captures students
    who arrive late or leave early, replacing the single-shot model that
    previous edge-based systems typically use.

  \item \textbf{A two-tier AI analytics layer.}  Deterministic trend
    classification for robustness combined with LLM-generated prose for
    interpretability---a pattern that prevents hallucinated structured
    values while preserving the human-readable output that academic
    advisers require.

  \item \textbf{A 5$\times$ pipeline runtime reduction.}  Batch-query
    optimisation reduced the at-risk explanation pipeline from 14--16
    minutes to approximately 2--3 minutes for 8 students, making the
    feature practical for page-triggered on-demand use.

  \item \textbf{A 22-phase iterative delivery within one semester.}
    The phased development model allowed early end-to-end integration
    (hardware $\rightarrow$ firmware $\rightarrow$ backend $\rightarrow$
    frontend) and frequent risk mitigation, with stub and mock modes
    enabling software progress independent of hardware availability.
\end{enumerate}

\subsection{Future Work}
\label{subsec:future}

Several extensions are identified for future development:

\begin{itemize}
  \item \textbf{Multi-room support.}  A shared backend with
    per-room MQTT namespacing, database partitioning by
    \texttt{room\_id}, and a room-selection UI for professors.

  \item \textbf{React Native mobile app.}  A companion app for
    professors to view live sensor data, acknowledge alerts, and start
    sessions from a smartphone.

  \item \textbf{IR illuminator.}  Adding an IR LED ring to the camera
    module to improve low-light face recognition accuracy from 71\% to
    near-normal-light levels.

  \item \textbf{HTTPS/WSS deployment.}  SSL termination in nginx using
    a university-issued LAN certificate, with no backend changes
    required.

  \item \textbf{Quantised GGUF model via llama.cpp.}  Replacing Ollama
    with a 4-bit quantised phi3-mini GGUF model to reduce inference time
    from 10--30~s to an estimated 3--8~s per student on the RPi~4B CPU.

  \item \textbf{Calibrated CO$_2$ sensor.}  Replacing the MQ-135 with
    an SCD30 or MH-Z19B NDIR sensor for accurate CO$_2$ monitoring and
    tighter air-quality thresholding.

  \item \textbf{Daily at-risk email digest.}  Scheduled morning emails
    to academic advisers summarising students whose attendance fell below
    threshold since the previous digest.

  \item \textbf{QR-code self-enrollment.}  Allowing students to
    upload their own enrollment photos via a QR-linked web page,
    reducing the manual enrollment burden on professors.

  \item \textbf{Offline SQLite fallback.}  A lightweight persistence
    layer that continues recording attendance when the PostgreSQL
    container is unreachable, with automatic sync on recovery.
\end{itemize}
